\section{Integration and subgroup congruences}\label{sec:congruences}

Fix $l\geq1$ and retain the group $\Gamma$, the tower, and the notation of
\Cref{thm:pro-characters}.  The purpose of this section is to extract arithmetic
information from the natural evaluation splitting.  Integration determines the
weights of its coordinates.  Integral cochains then force congruences between
weighted counts of finite sets with different actions.

For a finite groupoid $Y$, write
\[
 A_Y=\one^Y,\qquad
 \mathcal X_\Gamma(Y)=\Fun_{\cts}(B\Gamma,Y),\qquad
 a_\alpha=|\Aut_{\mathcal X_\Gamma(Y)}(\alpha)|.
\]
We use finite groupoid models and continuous functors throughout.
The groupoid $\mathcal X_\Gamma(Y)$ is finite up to equivalence: on each
component $BF$ there are finitely many homomorphisms $\Gamma\to F$, since
$\Gamma$ is finitely generated.  Its \emph{homotopy cardinality} is
\[
 |\mathcal X_\Gamma(Y)|_{\mathrm h}
   =\sum_{[\alpha]\in\pi_0\mathcal X_\Gamma(Y)}\frac1{a_\alpha}.
\]
All integrations below are first formed in $\Sp_{T(l)}$ and then transported
to $\mathcal C_{\Q}$.  We use their composition, base-change, and module
identities, and their agreement with the sum transfer for finite coverings;
these follow from finite-groupoid ambidexterity
\cite[Theorem~5.3.1, Propositions~3.1.10, 3.1.14, and 3.3.1]{CSY}.
In particular, we do not form new cochain objects after rationalizing morphisms.

\subsection{Absolute and relative integration weights}

Choose a common sufficiently advanced stage $L$ for the finite groupoids under
consideration and put $D=D_L$.  Under \Cref{thm:pro-characters} we have
\begin{equation}\label{eq:groupoid-coordinates}
 D\otimes A_Y\simeq
 \prod_{[\alpha]\in\pi_0\mathcal X_\Gamma(Y)}D.
\end{equation}
The coordinate maps are algebra maps.  We denote evaluation at $\alpha$ by
$\chi_\alpha$, and write $e_\alpha$ for the corresponding coordinate idempotent.
Scalar values lie in the ring
$R_D=\Hom_{\mathcal C_{\Q}}(\one,D)$. We also write $\epsilon_Y$
and $f_!$ for their $D$-linear extensions; tensor products between
extended cochain algebras are taken over $D$.

For characters indexed by maps from $\Z_p^l$, the same automorphism weights
as in the following proposition occur in Strickland's Frobenius character
formulas \cite[Section~10 and Theorem~10.2]{StricklandDuality}.
We derive them for $\Gamma$ directly from the diagonal transfer and the
evaluation splitting of \Cref{thm:pro-characters}.

\begin{proposition}[Integration weights]\label{prop:relative-weights}
In the coordinates \eqref{eq:groupoid-coordinates}, the $D$-linear
extension of integration $\epsilon_Y:A_Y\to\one$ is the map
\begin{equation}\label{eq:absolute-weights}
 (x_\alpha)_{[\alpha]}\longmapsto
       \sum_{[\alpha]}\frac{x_\alpha}{a_\alpha}.
\end{equation}
For a map $f:Y\to Z$, its pushforward has coordinate formula
\begin{equation}\label{eq:relative-weights}
 \bigl(f_!(x)\bigr)_\beta
   =\sum_{[\alpha]:[f\alpha]=[\beta]}
       \frac{a_\beta}{a_\alpha}\,x_\alpha.
\end{equation}
Consequently $\chi_\beta(f_!1)$ is the rational scalar
\begin{equation}\label{eq:relative-fiber-mass}
 \left|
   \mathcal X_\Gamma(Y)
       \mathop{\times}^{h}_{\mathcal X_\Gamma(Z)}\{\beta\}
 \right|_{\mathrm h}
\end{equation}
times the unit of $D$.
\end{proposition}

\begin{proof}
The diagonal $\Delta:Y\to Y\times Y$ is a finite covering: its homotopy fiber
at $(y,y')$ is the finite set $\operatorname{Iso}_Y(y,y')$.
By \Cref{lem:cover-evaluation}, its transferred unit evaluates at a pair
$(\alpha,\beta)$ as the number of isomorphisms between the two continuous
functors.  Thus
\[
 \Delta_!1=\sum_{[\alpha]}a_\alpha(e_\alpha\otimes e_\alpha)
\]
in the coordinates for $Y\times Y$.  Composition of integrations, applied to
$\operatorname{pr}_1\Delta=\id_Y$, gives
\[
 (\id\otimes\epsilon_Y)(\Delta_!1)=1.
\]
At the $\alpha$ coordinate this reads
$a_\alpha\epsilon_Y(e_\alpha)=1$.  Positive integers are invertible in
$\mathcal C_{\Q}$, proving \eqref{eq:absolute-weights}.

Naturality of the evaluation splitting identifies $f^*$ with pullback along
$\pi_0\mathcal X_\Gamma(Y)\to\pi_0\mathcal X_\Gamma(Z)$.
The module identity for $f_!$ therefore forces $f_!(e_\alpha)$ to be supported
at $[f\alpha]$.  If its value there is $c_\alpha$, composition with integration
to the point gives
\[
 \frac{c_\alpha}{a_{f\alpha}}=\frac1{a_\alpha}.
\]
This proves \eqref{eq:relative-weights}.

For completeness, the homotopy fiber in \eqref{eq:relative-fiber-mass} has
objects $(\alpha,\eta)$ with $\eta:f\alpha\xrightarrow{\sim}\beta$.
The contribution of a fixed isomorphism class $[\alpha]$ is the cardinality
of the action groupoid
\[
 \operatorname{Iso}(f\alpha,\beta)\mathbin{//}\Aut(\alpha),
\]
which is $a_\beta/a_\alpha$ if $f\alpha\simeq\beta$ and zero otherwise.
This computation includes automorphisms in the kernel of
$\Aut(\alpha)\to\Aut(f\alpha)$, so it applies to every map of finite
groupoids.
\end{proof}

\subsection{Integral scalar values}

The next argument concerns a cochain whose selected character value has been
computed as a rational scalar.  It uses invertibility in the original local
category to constrain that scalar.

\begin{lemma}[Scalar integrality]\label{lem:scalar-integrality}
Let $F$ be finite, let $z\in\pi_0A_F$, and suppose that an evaluation at an
advanced stage satisfies $\chi_\alpha(z)=r\cdot1_D$ for $r\in\Q$.
Then $r\in\Z_{(p)}$.  If the restriction of $z$ to the underlying point of
$BF$ is zero in $\pi_0\one$, then $r\in p\Z_{(p)}$.
\end{lemma}

\begin{proof}
We first show that $1+px$ is a unit in $\pi_0\one$ for every
$x\in\pi_0\one$.  Put $M=S/p$.  In the $p$-local stable category,
$M$-acyclic spectra have $p$ invertible and hence are rational.  They are
therefore $T(l)$-acyclic for $l>0$.  It follows that every $T(l)$-local
spectrum is $M$-local.

The cofiber triangle
\[
 S\xrightarrow{p}S\xrightarrow{i}M\xrightarrow{\partial}\Sigma S
\]
gives $p^2\id_M=0$, including when $p=2$.  Indeed,
$(p\id_M)i=ip=0$, so $p\id_M$ factors as $a\partial$ for a map
$a:\Sigma S\to M$.  Hence
\[
 p^2\id_M=p a\partial=a(\Sigma p)\partial=0.
\]
After smashing with $M$, the endomorphism $1+px$ of $\one$ thus has inverse
$1-px$.  Its cofiber is both $M$-local and $M$-acyclic, so it is zero.
This proves the asserted invertibility.

An element of $\pi_0A_F$ is an endomorphism of the constant unit in
$\Fun(BF,\Sp_{T(l)})$.  Evaluation at the unique object of $BF$ detects
equivalences of diagrams.  Consequently a cochain whose point restriction
is a unit in $\pi_0\one$ is itself a unit in $\pi_0A_F$.

Suppose $r\notin\Z_{(p)}$.  Then $r^{-1}\in p\Z_{(p)}$, and the
\emph{integral} cochain
\[
 u=1-r^{-1}z\in\pi_0A_F
\]
has point restriction of the form $1+px$.  It is a unit, whereas
$\chi_\alpha(u)=0$.  An algebra map cannot send a unit to zero in the
nonzero algebra $D$, a contradiction.  This proves $r\in\Z_{(p)}$.
If $z$ has zero point restriction and $r\notin p\Z_{(p)}$, then $r$ is a
unit of $\Z_{(p)}$; the same cochain $u$ now has point restriction exactly
$1$, giving the same contradiction.
\end{proof}

\begin{lemma}[Integrality of homomorphism counts]\label{lem:hom-integrality}
For every finite group $F$,
\begin{equation}\label{eq:hom-integrality}
 \frac{|\Hom_{\cts}(\Gamma,F)|}{|F|}\in\Z_{(p)}.
\end{equation}
\end{lemma}

\begin{proof}
The integral $\epsilon_{BF}(1)$ is an actual element of $\pi_0\one$.
By \Cref{prop:relative-weights}, its image in a sufficiently advanced
nonzero $D$ is the scalar
\[
 \sum_{[\alpha]\in\Hom_{\cts}(\Gamma,F)/F}
       \frac1{|C_F(\im\alpha)|}
   =\frac{|\Hom_{\cts}(\Gamma,F)|}{|F|},
\]
where the equality is orbit--stabilizer for conjugation by $F$.
Apply \Cref{lem:scalar-integrality} with the trivial group as its $F$.
\end{proof}

For the fixed source $\Z_p^l$, Ben-Moshe proves compatibility of
transchromatic characters with integration.  His cardinality formula on
$BF$ gives
$|\Hom_{\cts}(\Z_p^l,F)|/|F|\in\Z_{(p)}$
\cite[Theorems~5.7 and~5.8]{BenMoshe}.
Here the source $\Gamma$ is selected by the occurrence construction of
\Cref{sec:characters}.

For finitely generated discrete groups, Gordon and Rodriguez-Villegas
characterize integrality of these homomorphism-count ratios for all finite
targets by infinite abelianization \cite[Corollary~3.3]{GRV}.
Here the homomorphisms are continuous and the conclusion is $p$-integrality;
cyclic targets will give the first $\Z_p$ quotient in \Cref{sec:profinite}.

\subsection{Labeled sets and open subgroups}

\begin{theorem}[Congruences for open subgroups]\label{thm:subgroup-congruence}
For every open subgroup $U\le_o\Gamma$, every finite continuous $U$-set
$X$, and every $h\geq0$,
\begin{equation}\label{eq:subgroup-congruence}
 \sum_{\substack{V\le_o U\\{}[U:V]=p^h}}
       \bigl(|X^V|-|X|\bigr)
       \in p^{h+1}\Z_{(p)}.
\end{equation}
The sum is over all open subgroups of the indicated index.
\end{theorem}

\begin{proof}
We realize weighted counts of finite $U$-sets equivariantly labeled in
$X$ as character values of integral pushforwards. The logarithm of
their generating function isolates the transitive sets, whose
stabilizers give the open subgroups in \eqref{eq:subgroup-congruence}.

For $j\geq0$, let $\mathcal L_j(X)$ be the groupoid of continuous
$U$-sets $T$ of size $j$ equipped with a $U$-equivariant map $T\to X$.
Its arrows are the equivariant bijections preserving the labeling.  Since
$U$ is finitely generated, it has only finitely many homomorphisms to
$\Sigma_j$, so this groupoid is finite up to equivalence.  Set
\begin{equation}\label{eq:labeled-count}
 b_j(X)=|\mathcal L_j(X)|_{\mathrm h}
       =\sum_{[T\to X]}\frac1{|\Aut_U(T\to X)|}.
\end{equation}
Write $X_{\mathrm{triv}}$ for the same underlying set with trivial action.
We first prove
\begin{equation}\label{eq:labeled-integrality}
 b_j(X)\in\Z_{(p)},\qquad
 b_j(X)-b_j(X_{\mathrm{triv}})\in p\Z_{(p)}.
\end{equation}

Choose $N\trianglelefteq_o\Gamma$ contained in $U$ and in the kernel of
the action on $X$, and put $Q=\Gamma/N$ and $\overline U=U/N$.
The finite groupoid
\[
 Y_j(X)=X^j\mathbin{//}\Sigma_j
\]
classifies $j$-element sets labeled in $X$; $\overline U$ acts by changing
the labels.  Form its coinduction
\[
 Z_j(X)=\operatorname{Coind}_{\overline U}^{Q}Y_j(X).
\]
This is a finite $Q$-groupoid whose underlying groupoid is a product of
$[\Gamma:U]$ copies of $Y_j(X)$.  More explicitly, choosing representatives
$r_t$ for $t\in Q/\overline U$, with $r_{\overline U}=1$, gives this product
model; if
$a r_t=r_{at}u(a,t)$, then $a$ sends the $t$-coordinate to the
$at$-coordinate using the action of $u(a,t)\in\overline U$.
Let
\[
 \pi_j:E_j(X)=Z_j(X)\mathbin{//}Q\longrightarrow BQ,
 \qquad z_j(X)=(\pi_j)_!1\in\pi_0A_Q.
\]
These are a map of finite groupoids and an integral cochain.

Make the same construction for $X_{\mathrm{triv}}$, using the same $N$ and
$Q$.  For this fixed $j$, choose a common stage $L$ sufficiently advanced
for $Q$ and every component automorphism group of both total groupoids
$E_j(X)$ and $E_j(X_{\mathrm{triv}})$.  Then $D=D_L\ne0$, and inflation gives
equivalences
\[
 \Fun(BL,E_j(X))\simeq\mathcal X_\Gamma(E_j(X)),
 \qquad
 \Fun(BL,BQ)\simeq\mathcal X_\Gamma(BQ),
\]
and likewise for $X_{\mathrm{triv}}$.
In particular, the quotient homomorphism $q:\Gamma\to Q$ factors through
$L$.  Factorization gives essential surjectivity of these equivalences,
and surjectivity of $\Gamma\to L$ identifies the intertwining conditions
and gives full faithfulness.  The equivalences therefore identify homotopy
fibers over $q$.
The quotient $Q$ specifies the original action, while $L$ also captures
actions on the labeled objects.  Each coefficient uses finitely many total
groupoids, so such a common stage exists by \Cref{thm:pro-characters}.

We can now evaluate both cochains at the character indexed by $q$ over
this same $D$.  By \Cref{prop:relative-weights}, the value of $z_j(X)$ is
the homotopy cardinality of
\begin{equation}\label{eq:coinduced-fiber}
 \begin{aligned}
 \mathcal X_\Gamma(E_j(X))
       \mathop{\times}^{h}_{\mathcal X_\Gamma(BQ)}\{q\}
 &\simeq Z_j(X)^{h\Gamma}\\
 &\simeq Y_j(X)^{hU}
 \simeq\mathcal L_j(X).
 \end{aligned}
\end{equation}
Here $(-)^{h\Gamma}$ and $(-)^{hU}$ mean continuous equivariant objects,
including their coherence isomorphisms.  The first equivalence identifies
a lift of $q$, together with its specified identification over $BQ$, with
an equivariant object of the pulled-back bundle.
For the second, evaluate an equivariant object at the identity coset.
To construct an inverse, choose lifts $g_t\in\Gamma$ of the representatives
$r_t$, with $g_U=1$, and place a copy of the $U$-equivariant object in each
coordinate.
The coherence map for $\gamma\in\Gamma$ from coordinate $t$ to coordinate
$\gamma t$ is supplied by the $U$-equivariance map for
\[
 g_{\gamma t}^{-1}\gamma g_t\in U.
\]
Its image in $\overline U$ is $u(q\gamma,t)$, as required by the product
action above.
The cocycle identity proves coherence; the same construction on morphisms
shows that the two functors are inverse equivalences.
Continuity follows from continuity of the $U$-equivariance data and the
finiteness of $\Gamma/U$.
Finally, equivariance of a labeled $j$-element set is precisely a continuous
$U$-action on that set for which its map to $X$ is equivariant.
Thus $\chi_q(z_j(X))=b_j(X)\cdot1_D$.
The same argument gives
$\chi_q(z_j(X_{\mathrm{triv}}))=b_j(X_{\mathrm{triv}})\cdot1_D$.

After forgetting their $Q$-actions, the groupoids $Z_j(X)$ and
$Z_j(X_{\mathrm{triv}})$ are equivalent: their product factors have the same
underlying labeled-set groupoid.  Base change along $*\to BQ$ therefore
gives the integral equality
\[
 z_j(X)|_* =\epsilon_{Z_j(X)}(1)
           =\epsilon_{Z_j(X_{\mathrm{triv}})}(1)
           =z_j(X_{\mathrm{triv}})|_*
       \quad\text{in }\pi_0\one.
\]
Apply \Cref{lem:scalar-integrality} to $z_j(X)$ and to the difference of the
two classes, evaluated over the same $D$.  This proves
\eqref{eq:labeled-integrality}.

It remains to pass from all labeled sets to transitive ones.  Define the
formal power series
\[
 B_X(t)=\sum_{j\geq0}b_j(X)t^j\in1+t\Z_{(p)}[[t]].
\]
Every finite labeled $U$-set is a disjoint union of transitive labeled
$U$-sets.  If a transitive isomorphism type $a$ has size $m_a$ and
automorphism group of order $c_a$, then $k$ copies have automorphism group
of order $c_a^k k!$.  This orbit decomposition is the mechanism behind
Tamanoi's automorphism-weighted exponential formulas
\cite[Theorem~5-5 and equation~(5-7)]{Tamanoi2003}.
Consequently, degree by degree in $\Q[[t]]$,
\begin{equation}\label{eq:transitive-exponential}
 B_X(t)=\prod_a\sum_{k\geq0}\frac{t^{k m_a}}{c_a^k k!},
 \qquad
 \log B_X(t)=\sum_a\frac{t^{m_a}}{c_a}.
\end{equation}
There are finitely many types in each degree, which justifies these formal
products and sums.

A transitive continuous $U$-set has size $p^h$ for some $h$, since $U$ is
pro-$p$.  Pointing a transitive labeled set of size $p^h$ multiplies its
homotopy cardinality by $p^h$ and makes every automorphism trivial.
A pointed object is determined, up to unique isomorphism, by the stabilizer
$V\le_oU$ of its distinguished point and its label $x\in X^V$.
It follows that
\begin{equation}\label{eq:transitive-coefficient}
 [t^{p^h}]\log B_X(t)
   =\frac1{p^h}\sum_{\substack{V\le_oU\\{}[U:V]=p^h}}|X^V|.
\end{equation}
In particular this normalization counts actual subgroups, without passing
to conjugacy classes.

By \eqref{eq:labeled-integrality},
\[
 \frac{B_X(t)}{B_{X_{\mathrm{triv}}}(t)}=1+pA(t)
       \quad\text{for some }A(t)\in t\Z_{(p)}[[t]].
\]
Its formal logarithm lies in $p\Z_{(p)}[[t]]$, because
\[
 \log(1+pA(t))
   =\sum_{q\geq1}\frac{(-1)^{q+1}p^q}{q}A(t)^q,
 \qquad q-v_p(q)\geq1.
\]
Only finitely many terms contribute to any coefficient, and the valuation
inequality holds also at $p=2$.  Taking the coefficient of $t^{p^h}$,
using \eqref{eq:transitive-coefficient} for $X$ and $X_{\mathrm{triv}}$,
and multiplying by $p^h$ proves \eqref{eq:subgroup-congruence}.
\end{proof}
